% ---------------------------------------------------------------------------
% A starter entry. It exists so the pipeline has something to render and so you
% can see every supported construct in the browser. Delete it once you have
% real posts — or keep it and rewrite it, whichever is less work.
% ---------------------------------------------------------------------------

\title{Starting the Journal}
\date{2026-08-06 18:30}
\tags{life, hobbies}

\begin{document}

This is the first entry, which means it is mostly here to prove that the
machinery works. Everything below is written in LaTeX in
\texttt{blog/src/} and turned into this page by \texttt{build.py}.

Paragraphs are separated by a blank line, the way they are in any \LaTeX{}
document. Quotes and dashes get converted properly --- ``like this'' --- so I
can keep writing the way I already type.

\section{What goes here}

Food, travel, friends, hobbies, life. Those five tags sit at the top of the
page; clicking one narrows the feed and says so. Clicking it again clears the
filter.

\begin{itemize}
  \item \textbf{Food} --- what I cooked, what I ate, what went wrong.
  \item \textbf{Travel} --- where I went and whether it was worth the flight.
  \item \textbf{Friends} --- the people who make Atlanta feel like home.
  \item \textbf{Hobbies} --- tabletop games, side projects, whatever is current.
  \item \textbf{Life} --- everything that resists a neater label.
\end{itemize}

\begin{quote}
Time is steadily moving for all of us, which is the entire reason to write any
of this down.
\end{quote}

\section{Notes on the machinery}

Footnotes work and collect at the bottom of the entry.\footnote{This is a
footnote. It is not very interesting, but it renders.} So does inline math,
which I will inevitably need: the efficiency of a converter is
$\eta = P_{\mathrm{out}} / P_{\mathrm{in}}$, and display math gets its own line.

\[
  P_{\mathrm{loss}} = P_{\mathrm{cond}} + P_{\mathrm{sw}} + P_{\mathrm{core}}
\]

\sep

Images sit inline with a caption underneath.

\img{/assets/PESGM_group.jpg}{PES General Meeting.}

That is the whole system. Write a \texttt{.tex} file, run \texttt{python
build.py}, look at it locally, push when it reads right.

\end{document}
